\documentclass[11pt,reqno]{amsart}

\usepackage{amsmath,amssymb,amsthm}
\usepackage{booktabs}
\usepackage{hyperref}

\newtheorem{theorem}{Theorem}[section]
\newtheorem{proposition}[theorem]{Proposition}
\newtheorem{lemma}[theorem]{Lemma}
\newtheorem{corollary}[theorem]{Corollary}
\theoremstyle{definition}
\newtheorem{definition}[theorem]{Definition}
\theoremstyle{remark}
\newtheorem{remark}[theorem]{Remark}

\newcommand{\CLS}{\mathrm{CL}(S)}
\newcommand{\SN}{\mathrm{SN}}
\newcommand{\HNF}{\mathrm{HNF}}
\newcommand{\NF}{\mathrm{NF}}
\newcommand{\toS}{\to_S}
\newcommand{\redS}{\twoheadrightarrow_S}
\newcommand{\convS}{=_S}

\title[Short Proofs for Questions 7 and 31 in ``Dance of the Starlings'']{Short, Automaton-Free Proofs for Questions 7 and 31\\in ``Dance of the Starlings''}

\author{Guanquan Wen}
\thanks{Formalization and companion materials are archived at \url{https://github.com/Forage409/pure-s-combinator-universality-lean4} and \url{https://doi.org/10.5281/zenodo.24550796}. Multi-agent AI tools assisted in proof exploration and Lean~4 formalization; every proof in this note is short enough to be verified by hand and is independently checked in the standalone zero-import Lean~4.33.1 file \texttt{SIMPLIFIED\_Q31\_Q7\_BOOK\_PROOFS.lean} with zero \texttt{sorry}, zero \texttt{native\_decide}, and zero \texttt{Classical.choice}.}

\begin{document}

\begin{abstract}
We present short, self-contained, automaton-free proofs of two open questions posed by Barendregt, Endrullis, Klop, and Waldmann in \emph{Dance of the Starlings} (2017): Question~31 (whether $A A A$ and $A A A S$ are convertible in pure-$S$ combinatory logic $\CLS$) and Question~7 (the colour-change puzzle for strong normalization under single-leaf refinement). Both arguments share a single three-constructor inductive invariant $\mathcal{H}$ and its triple-application class $\mathcal{U} = \mathcal{H}\mathcal{H}\mathcal{H}$. For Question~31, this invariant combines with right-child persistence and confluence to eliminate our earlier $36$-state separating $S$-algebra and prove the stronger statement that $M S^j \neq_S M S^k$ for all $j \neq k$ whenever $M$ has no head normal form and has a compound right child; in particular, $\{A A A S^k \mid k \in \mathbb{N}\}$ yields infinitely many distinct conversion classes sharing the clocked B\"ohm tree $\bot$. For Question~7, the same invariant gives a two-step proof of non-termination for $R(S) = S A S A$, while $R(S S)$ terminates in a $6$-term reduction graph. This yields both a strictly alternating single-leaf refinement chain and an atomic leaf-splitting sequence ($S \mapsto S S$, adding one leaf per step) satisfying the exact characterization $T_n \in \SN \iff n \not\equiv 9 \pmod{10}$.
\end{abstract}

\maketitle

\section{Introduction and Preliminaries}\label{sec:intro}

Pure-$S$ combinatory logic $\CLS$ is the term rewriting system generated by a single constant $S$, variables, binary left-associative application, and the rewrite rule
\[
  S x y z \toS (x z)(y z).
\]
We write $\redS$ for the reflexive-transitive closure of $\toS$, $\convS$ for its symmetric-transitive closure ($S$-convertibility), and abbreviate
\[
  B := S S, \qquad A := S S S.
\]
The \emph{left-spine length} $\mathrm{sp}(M)$ is defined by $\mathrm{sp}(S) = \mathrm{sp}(x) = 0$ and $\mathrm{sp}(M N) = 1 + \mathrm{sp}(M)$. A closed term $M$ is in \emph{head normal form} ($\HNF$) if $\mathrm{sp}(M) \le 2$ (equivalently, $M$ has the shape $S$, $S x$, or $S x y$); $M$ \emph{has a head normal form} ($\mathrm{HasHNF}(M)$) if $M \redS H$ for some $H$ with $\mathrm{sp}(H) \le 2$. A term $M$ is \emph{strongly normalizing} ($M \in \SN$) if there is no infinite reduction sequence $M = M_0 \toS M_1 \toS \cdots$, and is in \emph{normal form} ($\NF$) if no $M'$ satisfies $M \toS M'$.

In \emph{Dance of the Starlings}~\cite{BarendregtEndrullisKlopWaldmann2017}, Barendregt, Endrullis, Klop, and Waldmann asked:
\begin{enumerate}
  \item \textbf{Question~31} (\cite[p.~93]{BarendregtEndrullisKlopWaldmann2017}): Are $A A A$ and $A A A S$ convertible in $\CLS$? Both terms have no head normal form and therefore have identical clocked B\"ohm trees $\bot$; a negative answer shows that clocked B\"ohm trees are not faithful for $S$-convertibility.
  \item \textbf{Question~7} (\cite[p.~81]{BarendregtEndrullisKlopWaldmann2017}): Colour a closed $S$-term green if it is strongly normalizing and red otherwise. Can one start from $S$ and repeatedly refine single $S$-leaves so that the colour alternates infinitely often between green and red?
\end{enumerate}

In an earlier monograph, we resolved Question~31 using a $36$-state $S$-algebra $\mathcal{A}_{36}$ and used Waldmann's $35$-state solvability algebra $\mathcal{A}_{\mathrm{HN}}$ inside the proof of Question~7. Following a suggestion of J\"org Endrullis to simplify these arguments as far as possible, we show here that neither finite algebra is needed for either result. Both theorems follow cleanly from a single three-constructor inductive set $\mathcal{H}$ (a streamlined form of Waldmann's $QQQ$ invariant~\cite{Waldmann2000}) together with elementary syntactic observations and standard confluence of $\CLS$.

\begin{lemma}[Church--Rosser for $\CLS$~\cite{Takahashi1995}]\label{lem:cr}
If $M \convS N$, then there exists a term $W$ such that $M \redS W$ and $N \redS W$.
\end{lemma}

In our accompanying Lean~4 file, Lemma~\ref{lem:cr} is proved from scratch in roughly forty lines via Takahashi's complete development method~\cite{Takahashi1995} without importing external libraries.

\section{A Shared Three-Constructor Invariant}\label{sec:invariant}

\begin{definition}\label{def:heavy}
Let $\mathcal{H}$ be the smallest set of closed pure-$S$ terms satisfying:
\begin{enumerate}
  \item[(H1)] $A \in \mathcal{H}$;
  \item[(H2)] if $x \in \mathcal{H}$, then $S x \in \mathcal{H}$;
  \item[(H3)] if $x, y \in \mathcal{H}$, then $x y \in \mathcal{H}$.
\end{enumerate}
Define the triple-application class $\mathcal{U} := \{\, p q r \mid p, q, r \in \mathcal{H} \,\}$.
\end{definition}

\begin{lemma}[Inversion properties of $\mathcal{H}$]\label{lem:heavy-inv}
\begin{enumerate}
  \item $S \notin \mathcal{H}$ and $\mathrm{sp}(t) \ge 1$ for every $t \in \mathcal{H}$.
  \item If $S x \in \mathcal{H}$, then $x \in \mathcal{H}$.
  \item If $S x y \in \mathcal{H}$, then either $(x, y) = (S, S)$ or $x, y \in \mathcal{H}$.
\end{enumerate}
\end{lemma}

\begin{proof}
Clause~(1) is immediate by inspection of (H1)--(H3), since $A = (S S) S$, $S x$, and $x y$ are all applications. For~(2), $S x$ cannot arise from (H1) because its left child is $S \neq S S$, and cannot arise from (H3) because $S \notin \mathcal{H}$; thus it must arise from (H2), giving $x \in \mathcal{H}$. For~(3), $S x y = (S x) y$ cannot arise from (H2) because $S x \neq S$. If it arises from (H1), then $(x, y) = (S, S)$; if it arises from (H3), then $S x \in \mathcal{H}$ and $y \in \mathcal{H}$, whence $x \in \mathcal{H}$ by~(2).
\end{proof}

\begin{proposition}[Closure and divergence of $\mathcal{U}$]\label{prop:U}
\begin{enumerate}
  \item If $t \in \mathcal{H}$ and $t \toS t'$, then $t' \in \mathcal{H}$.
  \item If $t \in \mathcal{U}$ and $t \toS t'$, then $t' \in \mathcal{U}$.
  \item Every $t \in \mathcal{U}$ satisfies $\mathrm{sp}(t) \ge 3$, $\neg \mathrm{HasHNF}(t)$, and $t \notin \SN$.
\end{enumerate}
\end{proposition}

\begin{proof}
(1) Proceed by induction on $t \in \mathcal{H}$. For (H1), $A = S S S$ is a normal form, so no step occurs. For (H2), a step from $S x$ must reduce $x \toS x'$, keeping $S x' \in \mathcal{H}$ by the induction hypothesis. For (H3), let $t = u v$ with $u, v \in \mathcal{H}$. Reductions inside $u$ or $v$ preserve $\mathcal{H}$ by the induction hypothesis. If a root redex fires, then $u = S x y$ and $t' = (x v)(y v)$. By Lemma~\ref{lem:heavy-inv}(3), either $x = y = S$, giving $t' = (S v)(S v) \in \mathcal{H}$ by (H2) and (H3), or $x, y \in \mathcal{H}$, giving $t' = (x v)(y v) \in \mathcal{H}$ by (H3).

(2) Let $t = p q r$ with $p, q, r \in \mathcal{H}$, and suppose $t \toS t'$. If the step occurs inside $p$, $q$, or $r$, then $t' \in \mathcal{U}$ by~(1). If a root redex fires at $p q$, then $p = S x y$ and $p q \toS (x q)(y q)$; by Lemma~\ref{lem:heavy-inv}(3), both $x q$ and $y q$ lie in $\mathcal{H}$ (taking $S q \in \mathcal{H}$ when $x = y = S$), so $t' = (x q)(y q) r \in \mathcal{U}$. Finally, if a root redex fires at the top of $p q r$, then $p q = S x y$, so $p = S x$ and $q = y$. By Lemma~\ref{lem:heavy-inv}(2), $x \in \mathcal{H}$; hence $t' = (x r)(q r) = p' q' r' \in \mathcal{U}$ with $p' = x$, $q' = r$, and $r' = q r \in \mathcal{H}$.

(3) For any $t = p q r \in \mathcal{U}$, Lemma~\ref{lem:heavy-inv}(1) gives $\mathrm{sp}(p) \ge 1$, so $\mathrm{sp}(t) = \mathrm{sp}(p) + 2 \ge 3$. By~(2), every reduct $t'$ of $t$ lies in $\mathcal{U}$ and therefore satisfies $\mathrm{sp}(t') \ge 3 > 2$, proving $\neg \mathrm{HasHNF}(t)$. Moreover, every $t \in \mathcal{U}$ is closed with $\mathrm{sp}(t) \ge 3$, so its leftmost leaf is $S$ followed by at least three arguments on the left spine; thus $t$ contains a head redex $S x y z$ and admits a step $t \toS t' \in \mathcal{U}$. Repeating this step produces an infinite reduction sequence, so $t \notin \SN$.
\end{proof}

\section{Question 31: Pairwise Non-Convertibility of $\{A A A S^k\}_{k \ge 0}$}\label{sec:q31}

Call a term $M$ \emph{compound} if $M = u v$ for some terms $u, v$, and \emph{right-compound} if $M = u v$ with $v$ compound. For any term $M$ and $k \in \mathbb{N}$, define $M S^0 := M$ and $M S^{k+1} := (M S^k) S$. Let $\tau(M)$ count the number of outer trailing applications to $S$:
\[
  \tau(M S) := 1 + \tau(M), \qquad \tau(M) := 0 \text{ if } M \text{ is not of the form } u S.
\]
In particular, if $M$ is right-compound, then $\tau(M) = 0$ and $\tau(M S^k) = k$ for every $k \in \mathbb{N}$.

\begin{lemma}[Right-child persistence]\label{lem:right-app}
If $M$ is right-compound and $M \redS M'$, then $M'$ is right-compound.
\end{lemma}

\begin{proof}
First observe that if $v$ is compound and $v \toS v'$, then $v'$ is compound (every constructor of $\toS$ yields an application). Now let $M = u v$ with $v$ compound, and suppose $M \toS M'$. If the step is a root redex $(S x y) z \toS (x z)(y z)$, the new right child $y z$ is compound. If the step reduces $u \toS u'$, the right child $v$ is unchanged. If the step reduces $v \toS v'$, the new right child $v'$ remains compound. Induction on the length of $M \redS M'$ completes the proof.
\end{proof}

\begin{lemma}[Trailing-$S$ freeze]\label{lem:freeze}
Suppose $\neg \mathrm{HasHNF}(M)$. Then for every $k \in \mathbb{N}$ and every term $W$ with $M S^k \redS W$, there exists $M'$ such that $M \redS M'$ and $W = M' S^k$.
\end{lemma}

\begin{proof}
We proceed by induction on $k$; the case $k = 0$ is trivial. Assume the statement holds for $k$, and consider any reduction $M S^{k+1} \redS W$. We show by inner induction on the number of reduction steps that every term reachable from $M' S^{k+1}$ (where $M \redS M'$) has the form $M'' S^{k+1}$ with $M \redS M''$. Take a single step $M' S^{k+1} = (M' S^k) S \toS V$:
\begin{itemize}
  \item A step inside the right child $S$ is impossible because $S$ is irreducible.
  \item A root step requires $M' S^k = S x y$ for some $x, y$. If $k = 0$, then $M' = S x y$; if $k = 1$, then $M' = S x$; if $k = 2$, then $M' = S$; and if $k \ge 3$, then $M' S^{k-2} = S$, which is impossible since $M' S^{k-2}$ is an application. In the cases $k \in \{0, 1, 2\}$, we obtain $M \redS M'$ with $\mathrm{sp}(M') \le 2$, contradicting $\neg \mathrm{HasHNF}(M)$.
  \item Hence the step must occur in the left child: $M' S^k \toS U$ with $V = U S$. Since $M S^k \redS M' S^k \toS U$, the outer induction hypothesis yields $M''$ with $M \redS M''$ and $U = M'' S^k$, so $V = M'' S^{k+1}$.
\end{itemize}
This closes both inductions.
\end{proof}

\begin{theorem}[General right-compound separation]\label{thm:general-q31}
Let $M$ be any right-compound term with $\neg \mathrm{HasHNF}(M)$. Then for all $j, k \in \mathbb{N}$ with $j \neq k$,
\[
  M S^j \neq_S M S^k.
\]
\end{theorem}

\begin{proof}
Suppose toward a contradiction that $M S^j \convS M S^k$. By Lemma~\ref{lem:cr}, there exists a common reduct $W$ with $M S^j \redS W$ and $M S^k \redS W$. Applying Lemma~\ref{lem:freeze} twice yields terms $M_1, M_2$ such that $M \redS M_1$, $M \redS M_2$, and
\[
  M_1 S^j = W = M_2 S^k.
\]
By Lemma~\ref{lem:right-app}, both $M_1$ and $M_2$ are right-compound, so $\tau(M_1) = \tau(M_2) = 0$. Taking $\tau$ on both sides gives $j = \tau(M_1 S^j) = \tau(W) = \tau(M_2 S^k) = k$, contradicting $j \neq k$.
\end{proof}

\begin{corollary}[Resolution of Question~31]\label{cor:q31}
For all $j \neq k \in \mathbb{N}$, $A A A S^j \neq_S A A A S^k$. In particular, $A A A \neq_S A A A S$.
\end{corollary}

\begin{proof}
Since $A \in \mathcal{H}$, we have $A A A \in \mathcal{U}$, so $\neg \mathrm{HasHNF}(A A A)$ by Proposition~\ref{prop:U}(3). The right child of $A A A = (A A) A$ is $A = (S S) S$, which is compound. Apply Theorem~\ref{thm:general-q31}.
\end{proof}

\begin{remark}
By Lemma~\ref{lem:freeze}, every $A A A S^k$ also has no head normal form (any reduct $M' S^k$ with $\mathrm{sp}(M' S^k) \le 2$ would satisfy $\mathrm{sp}(M') \le 2$, contradicting $\neg \mathrm{HasHNF}(A A A)$). Consequently, all infinitely many terms $\{A A A S^k \mid k \in \mathbb{N}\}$ share the clocked B\"ohm tree $\bot$, yet lie in pairwise distinct $S$-conversion classes. Exhaustive constraint-propagation search also verifies that no finite $S$-algebra of size $\le 5$ separates $A A A$ from $A A A S$, underscoring why the syntactic argument of Theorem~\ref{thm:general-q31} is substantially simpler than a finite-algebra separation such as our earlier $36$-state table $\mathcal{A}_{36}$.
\end{remark}

\section{Question 7: Automaton-Free Colour-Change Chains}\label{sec:q7}

Define the single-hole $8$-leaf gadget
\[
  R(x) := x A S A \qquad (A = S S S).
\]

\begin{lemma}[The red/green switch $R$]\label{lem:gadget}
$R(S) \notin \SN$ and $R(S S) \in \SN$.
\end{lemma}

\begin{proof}
For $R(S) = S A S A$, two reduction steps yield
\[
  S A S A \toS (A A)(S A) = ((S S S) A)(S A) \toS ((S A)(S A))(S A).
\]
Since $S A \in \mathcal{H}$ by (H1) and (H2), the term $((S A)(S A))(S A)$ belongs to $\mathcal{U}$. By Proposition~\ref{prop:U}(3), it is not strongly normalizing, and hence $R(S) \notin \SN$.

For $R(S S) = B A S A$ (recall $B = S S$), the terms $B$, $A$, $S A$, $B A$, and $B B$ are all normal forms. Direct inspection shows that the complete multiway reduction graph of $R(B)$ consists of exactly six terms:
\begin{align*}
  g_0 &= ((B A) S) A \toS g_1 := (B (A S)) A, \\
  g_1 &\toS g_2 := (S A)((A S) A) \qquad \text{and} \qquad g_1 \toS g_3 := (B (B B)) A, \\
  g_2 &\toS g_4 := (S A)((B B) A), \qquad g_3 \toS g_4, \\
  g_4 &\toS g_5 := (S A)((S A)(B A)) \in \NF.
\end{align*}
Since every reduction path from $g_0$ terminates at $g_5$ in at most $4$ steps, $R(S S) \in \SN$.
\end{proof}

We also use two elementary closure properties of $\SN$:
\begin{enumerate}
  \item[(P1)] If $u, v \in \SN$, then $S u v \in \SN$ (no root redex can ever form, so reductions occur independently inside $u$ and $v$).
  \item[(P2)] If $u v \in \SN$, then $u \in \SN$ and $v \in \SN$.
\end{enumerate}

Write $M \prec N$ if $N$ is obtained from $M$ by replacing a single $S$-leaf with a closed compound $S$-term, and write $M \prec_1 N$ for an \emph{atomic leaf split}, where a single $S$-leaf is replaced by $S S$ (so $|N| = |M| + 1$, where $|M|$ denotes the number of $S$-leaves of $M$).

Define the right-nested spine $\Phi_k(t)$ recursively by $\Phi_0(t) := t$ and $\Phi_{k+1}(t) := S\,R(S S)\,\Phi_k(t)$. Notice that $\Phi_{k+1}(S) = \Phi_k(S\,R(S S)\,S)$.

\begin{theorem}[Single-leaf refinement alternation for Question~7]\label{thm:q7-flip}
Define $F_{2k} := \Phi_k(S)$ and $F_{2k+1} := \Phi_k(S\,R(S)\,S)$ for all $k \in \mathbb{N}$. Then:
\begin{enumerate}
  \item $F_0 = S$, and $F_n \prec F_{n+1}$ for every $n \in \mathbb{N}$;
  \item $F_{2k} \in \SN$ (green) and $F_{2k+1} \notin \SN$ (red) for every $k \in \mathbb{N}$.
\end{enumerate}
\end{theorem}

\begin{proof}
At even indices, $F_{2k} = \Phi_k(S) \prec \Phi_k(S\,R(S)\,S) = F_{2k+1}$ replaces the rightmost leaf $S$ by $S\,R(S)\,S$. At odd indices, $F_{2k+1} = \Phi_k(S\,R(S)\,S) \prec \Phi_k(S\,R(S S)\,S) = \Phi_{k+1}(S) = F_{2k+2}$ replaces the head leaf $S$ of $R(S)$ by $S S$. By Lemma~\ref{lem:gadget} and induction on $k$ using (P1), $\Phi_k(S) \in \SN$ because $R(S S) \in \SN$ and $S \in \SN$. Conversely, $R(S)$ is a subterm of $F_{2k+1}$, so $F_{2k+1} \notin \SN$ by Lemma~\ref{lem:gadget} and (P2).
\end{proof}

Even under the strictest reading of Question~7---where every step is an atomic leaf split $S \mapsto S S$ increasing the leaf count by exactly $1$---the exact colour sequence is elementary. We define the $10$-term atomic bridge $P_0, \dots, P_9$ in Table~\ref{tab:atomic-bridge}.

\begin{table}[h]
\centering
\caption{The $10$-term atomic leaf-splitting bridge $P_0 \prec_1 \cdots \prec_1 P_9$.}\label{tab:atomic-bridge}
\begin{tabular}{rll}
\toprule
$r$ & $P_r$ & Strong normalization witness \\
\midrule
$0$ & $S$ & $\NF$ \\
$1$ & $S S$ & $\NF$ \\
$2$ & $S S S$ & $\NF$ \\
$3$ & $S (S S) S$ & $\NF$ \\
$4$ & $S A S$ & $\NF$ \\
$5$ & $S (A S) S$ & $A S \toS B B \in \NF$ ($2$ states) \\
$6$ & $S (A B) S$ & $A B \toS (S B)(S B) \in \NF$ ($2$ states) \\
$7$ & $S (A A) S$ & $A A \toS (S A)(S A) \in \NF$ ($2$ states) \\
$8$ & $S (S B S A) S$ & $4$-state linear chain to $\NF$ \\
$9$ & $S (S A S A) S = S\,R(S)\,S$ & non-$\SN$ (Lemma~\ref{lem:gadget}) \\
\bottomrule
\end{tabular}
\end{table}

\begin{theorem}[Exact colour pattern under atomic leaf splitting]\label{thm:q7-atomic}
With $P_0, \dots, P_9$ defined as in Table~\ref{tab:atomic-bridge}, set $T_{10k + r} := \Phi_k(P_r)$ for $k \in \mathbb{N}$ and $0 \le r \le 9$. Then $T_0 = S$, $T_n \prec_1 T_{n+1}$ and $|T_n| = n + 1$ for every $n \in \mathbb{N}$, and
\[
  T_n \in \SN \iff n \not\equiv 9 \pmod{10}.
\]
In particular, every step $10k + 8 \to 10k + 9$ switches from green to red, and the immediately following step $10k + 9 \to 10k + 10$ switches back from red to green.
\end{theorem}

\begin{proof}
Each transition $P_r \prec_1 P_{r+1}$ ($0 \le r \le 8$) and $P_9 = S\,R(S)\,S \prec_1 S\,R(S S)\,S$ replaces a single $S$-leaf by $S S$: for $r = 4 \to 5$, the leftmost leaf inside $A = (S S) S$ splits into $((S S) S) S = A S$; for $r = 7 \to 8$, the middle $S$ inside the left $A$ of $A A = (S S S) A$ splits to give $(S (S S) S) A$; for $r = 8 \to 9$, the first leaf of that $S S$ splits to give $S A S A$; and for $r = 9 \to 10$, the head $S$ of $S A S A$ splits to give $R(S S)$, matching $\Phi_{k+1}(S) = \Phi_k(S\,R(S S)\,S)$.

For any normal form $x$, $A x = S S S x$ has only the root redex $A x \toS (S x)(S x) \in \NF$, establishing $P_5, P_6, P_7 \in \SN$ via (P1). For $P_8$, the subterm $S B S A$ has a unique, unbranched $3$-step reduction to normal form:
\[
  S B S A \toS (B A)(S A) \toS (S (S A))(A (S A)) \toS (S (S A))\bigl((S (S A))(S (S A))\bigr) \in \NF,
\]
so $P_8 \in \SN$ by (P1). Thus $T_{10k + r} = \Phi_k(P_r) \in \SN$ for all $0 \le r \le 8$ by (P1) and Lemma~\ref{lem:gadget}, while $T_{10k + 9} \notin \SN$ by (P2) and Lemma~\ref{lem:gadget}.
\end{proof}

\begin{remark}[One-step rescue classification]\label{rem:rescue}
The colour-change chains above rely only on a single rescue transition $R(S) \mapsto R(S S)$ and need no finite tree automaton. By contrast, if one asks for a \emph{universal classification} of all closed pure-$S$ plugs $P$ capable of rescuing \emph{some} non-$\SN$ context $C[S]$ (meaning $C[S] \notin \SN$ while $C[P] \in \SN$), then Waldmann's minimal $39$-state algebra $\mathcal{A}_{\SN}$~\cite{Waldmann2000} (whose contextual preorder has $39$ distinct classes) does become essential: as formalized in \texttt{formal\_algebra/question7\_flipflop.lean} of our repository, the only closed pure-$S$ terms $P$ with $S \not\le_{\SN} P$ are syntactically $P = S S$ and $P = S(S S S)$, both of which are realized.
\end{remark}

\section{Formal Verification}\label{sec:lean}

Every definition, lemma, and theorem in Sections~\ref{sec:intro}--\ref{sec:q7} is formalized in the single self-contained Lean~4.33.1 file \texttt{SIMPLIFIED\_Q31\_Q7\_BOOK\_PROOFS.lean} (zero external imports, zero finite-algebra tables). Running \texttt{lean SIMPLIFIED\_Q31\_Q7\_BOOK\_PROOFS.lean} verifies the following theorems with axiom footprint \texttt{[propext, Quot.sound]} only (no \texttt{sorry}, no \texttt{native\_decide}, and no \texttt{Classical.choice}):
\begin{itemize}
  \item \texttt{SimplifiedQ31Q7.church\_rosser} (Lemma~\ref{lem:cr}),
  \item \texttt{SimplifiedQ31Q7.question\_31\_infinite\_family} and \texttt{question\_31\_aaa\_ne\_aaas} (Theorem~\ref{thm:general-q31} and Corollary~\ref{cor:q31}),
  \item \texttt{SimplifiedQ31Q7.question\_7\_infinite\_flipflop} (Theorem~\ref{thm:q7-flip}),
  \item \texttt{SimplifiedQ31Q7.question\_7\_atomic\_leaf\_splits} (Theorem~\ref{thm:q7-atomic}).
\end{itemize}

\begin{thebibliography}{9}

\bibitem{Barendregt1984}
H.~P. Barendregt, \emph{The Lambda Calculus: Its Syntax and Semantics}, revised ed., Studies in Logic and the Foundations of Mathematics, vol.~103, North-Holland, Amsterdam, 1984.

\bibitem{BarendregtEndrullisKlopWaldmann2017}
H.~Barendregt, J.~Endrullis, J.~W. Klop, and J.~Waldmann, \emph{Dance of the Starlings}, in \emph{Raymond Smullyan on Self Reference} (eds.\ M.~Fitting and B.~Rayman), Outstanding Contributions to Logic, vol.~14, Springer, Cham, 2017, pp.~67--111.

\bibitem{Takahashi1995}
M.~Takahashi, \emph{Parallel reductions in $\lambda$-calculus}, Information and Computation \textbf{118} (1995), no.~1, 120--127.

\bibitem{Waldmann2000}
J.~Waldmann, \emph{The combinator $S$}, Information and Computation \textbf{159} (2000), no.~1--2, 2--21.

\end{thebibliography}

\end{document}
